\chapter{Introducción general}
\label{chap:intro}

\section{Motivación}

Las crisis financieras de las últimas dos décadas han dejado una regularidad difícil de ignorar: cuando un sistema bancario frágil coincide con un deterioro agudo de las perspectivas de crecimiento, el costo de financiamiento de los gobiernos se dispara de forma más que proporcional a lo que cualquiera de los dos factores explicaría por separado. La literatura académica, sin embargo, ha tratado estos dos canales —la fragilidad bancaria sistémica y el riesgo a la baja del crecimiento— de manera predominantemente independiente. Por un lado, la literatura del nexo banca-soberano ha documentado extensamente cómo la fragilidad del sistema bancario se transmite al riesgo soberano a través de pasivos contingentes y rescates fiscales \citep{FarhiTirole2018,AcharyaDrechslerSchnabl2014}. Por otro, la literatura de \textit{Growth-at-Risk} ha caracterizado con precisión creciente cómo las condiciones financieras determinan de forma asimétrica la cola izquierda de la distribución del crecimiento \citep{ABG2019}. Ninguna de las dos tradiciones ha articulado sistemáticamente ambos canales como codeterminantes \textit{conjuntos y no lineales} del riesgo soberano.

Esta tesis nace de la conjetura de que esa articulación no es un simple ejercicio de agregar un regresor: la fragilidad bancaria y el riesgo de cola del crecimiento se \textbf{complementan}, no se suman, en su transmisión al spread soberano. La fragilidad del sistema bancario se comporta como un pasivo contingente cuyo costo esperado para el fisco depende críticamente del estado del ciclo económico; cuando el crecimiento ya está en su escenario adverso, el mismo nivel de fragilidad bancaria implica un rescate fiscal proporcionalmente más costoso, y por tanto una prima de riesgo soberano desproporcionadamente mayor.

\section{Dos aristas de una misma investigación}

Esta conjetura se investiga aquí desde dos ángulos deliberadamente distintos, que constituyen los dos capítulos centrales de la tesis.

El \textbf{Capítulo~\ref{chap:paper2}} aborda la arista \textbf{empírica}: estima directamente, sobre un único panel de trece economías emergentes con la fragilidad bancaria construida desde la terminal Bloomberg, el término de interacción entre la fragilidad bancaria sistémica —medida mediante la métrica \textit{Joint Loss} ($JLoss$) de \citet{Chari2024}— y el riesgo a la baja del crecimiento —medido mediante \textit{Growth-at-Risk} ($GaR$), siguiendo el marco de \textit{vulnerable growth} de \citet{ABG2019}—, sobre el spread soberano medido, siguiendo a \citet{Chari2024}, por el índice EMBI Global Diversified de J.P.\ Morgan. El capítulo establece con solidez los canales de \textbf{nivel} —la fragilidad bancaria y el riesgo de cola del crecimiento elevan cada uno el spread soberano— y muestra que su \textbf{complementariedad}, aunque del signo predicho y corroborada por un modelo de umbral, tiene una significancia estadística \textbf{condicional}: se sostiene sobre el núcleo de economías emergentes de financiamiento externo y se diluye en las economías con mercados de deuda local profundos. Somete el resultado a una batería exigente de robustez e identificación causal.

El \textbf{Capítulo~\ref{chap:paper1}} aborda la arista \textbf{teórica}, de organización industrial bancaria: desarrolla un modelo formal de competencia bancaria à la Cournot en el que la estructura de mercado —la intensidad de competencia entre bancos— es el parámetro profundo que gobierna toda la cadena, desde la fragilidad sistémica hasta el riesgo soberano. El modelo no solo deriva la complementariedad estimada en el Capítulo~\ref{chap:paper2} desde primeros principios, sino que genera una predicción adicional y distintiva frente a la literatura estándar de organización industrial bancaria \citep{MMR2010}: que esa complementariedad se \textbf{amplifica} en sistemas bancarios más concentrados. Esta predicción —la triple interacción entre fragilidad, riesgo de cola y concentración bancaria— se pone a prueba con datos reales al cierre del mismo capítulo, sin resultado concluyente: sobre el panel homogéneo, el coeficiente de amplificación no puede identificarse con potencia —su punto estimado no tiene signo estable entre proxies de concentración y su intervalo de confianza robusto cruza el cero con todos ellos—, y ese resultado se mantiene con una serie de concentración a frecuencia trimestral construida para este fin. El capítulo discute las explicaciones candidatas de esa discrepancia y la agenda de datos que permitiría dirimirla.

La decisión de presentar dos capítulos en lugar de fusionarlos en un único artículo obedece a una razón de fondo, no solo de extensión: son dos objetos de conocimiento distintos que se validan con estándares distintos. El Capítulo~\ref{chap:paper2} debe convencer con evidencia estadística —significancia, robustez, identificación causal—; el Capítulo~\ref{chap:paper1}, con rigor deductivo —existencia y unicidad del equilibrio, signos de estática comparativa derivados de supuestos explícitos—. Presentarlos por separado permite que cada uno satisfaga el estándar que le corresponde, mientras esta introducción y la discusión general del Capítulo~\ref{chap:discusion} muestran por qué, juntos, son más que la suma de sus partes: el modelo teórico explica por qué debería existir la complementariedad que el capítulo empírico documenta, y el capítulo empírico —al construir el mismo panel de datos que alimenta la prueba de la predicción distintiva del modelo (H4a y H4b, Sección~\ref{sec:datos_reales})— provee el terreno común donde ambas aristas se encuentran.

\section{Contribución y hoja de ruta}

En conjunto, esta tesis aporta: (i) evidencia empírica —sobre un único panel de trece economías emergentes con la fragilidad bancaria construida desde Bloomberg— de que la fragilidad bancaria sistémica y el riesgo a la baja del crecimiento elevan cada uno el spread soberano, y de que su \textbf{complementariedad} aparece en los precios de forma condicional: significativa en las economías más expuestas al financiamiento externo, no en el conjunto del panel; (ii) un modelo estructural de organización industrial bancaria que deriva esa complementariedad desde la competencia entre bancos, y que genera una predicción adicional —la amplificación por concentración— ausente de la literatura de organización industrial bancaria estándar; (iii) una prueba empírica de esa predicción distintiva sobre el mismo panel homogéneo, cuyo resultado no es concluyente y se reporta como tal, delimitando así qué parte del mecanismo estructural está hoy respaldada por los datos y qué parte permanece como hipótesis; y (iv) un \textit{pipeline} computacional íntegramente reproducible —documentado en el repositorio del proyecto, con una cadena de control de versiones explícita— para la construcción de ambas métricas y la estimación de ambos capítulos.

El resto de la tesis se organiza así: el Capítulo~\ref{chap:paper2} desarrolla la arista empírica; el Capítulo~\ref{chap:paper1}, la arista teórica y su prueba con datos reales; el Capítulo~\ref{chap:discusion} sintetiza ambas aristas, discute sus implicancias de política conjuntas y traza la agenda de investigación futura. El Anexo~\ref{chap:anexoA} contiene las demostraciones matemáticas completas del modelo teórico y el Anexo~\ref{chap:anexoB} tabula la procedencia de todas las series de datos.
